\section{Experiments}
\label{sec:experiments}
We evaluate the proposed LLM Performance Predictor (LPP) framework in the context of cost-aware selective escalation for human–AI content moderation. From a multi-agent systems perspective, our experiments study the coordination dynamics between autonomous LLM agents and human reviewers, where the LPP meta-model serves as a decision-making intermediary that allocates tasks under cost constraints.
\subsection{Experimental Setup}
\label{subsec:experimental_setup}
We evaluate several LLMs: proprietary (gpt-4o-mini, gpt-4o, gpt-4.1-mini, gemini-2.5-flash-lite, gemini-2.0-flash-lite, gemini-2.0-flash-001) and open-source (QWEN3-14B, QWEN3-32B, LLAMA32-11B).
\footnote{We considered including safety-specialized models such as Llama Guard~\cite{Inan2023LlamaGL}, but defer discussion of these to Section~\ref{subsec:limitations}.}
Following Section~\ref{sec:methods}, we instantiate LPPs using ridge regression trained on heterogeneous feature families derived from intrinsic model signals (e.g., token-level probabilities, self-consistency markers) and post-hoc signals (e.g., calibration-based scores). A summary of feature groups is provided in Table~\ref{tab:lpps}. Hyperparameters for baselines and the meta-model are tuned via nested cross-validation, with stratified sampling. To address class imbalance, we apply down-sampling as described in Sec.~\ref{subsec:training}, and evaluate final performance on a held-out test set using the optimized threshold.
The meta-model is trained with Ridge Regression, chosen for its robustness in high-dimensional feature spaces.


\myparagraph{Additional Experimental Details}
We use deterministic decoding parameters (\texttt{temperature}=0, \texttt{top-p}=1) with a fixed random seed (\texttt{random\_state}=42) for data splitting and cross-validation reproducibility. For each generated token, we request log-probabilities of the \textit{top-20} candidate tokens. From these we compute two feature families:
\begin{enumerate}
    \item \textbf{Unfiltered top-5 features}: standard uncertainty metrics computed over the five most probable tokens.
    \item \textbf{Filtered features}: log-probabilities restricted to the four valid class labels (\texttt{0-3}).
\end{enumerate}
This dual representation captures both the model’s general uncertainty landscape and its calibrated confidence over the task’s decision space. CoT perplexity is computed with natural logarithms. Cross-validation uses \texttt{StratifiedKFold} with $K=3$ splits.

\subsection{Datasets}
We use two complementary datasets:  
(i) the \textbf{OpenAI Moderation dataset}~\cite{markov2023holisticapproachundesiredcontent}, which contains 1,680 English text samples annotated across multiple moderation categories (hate, self-harm, sexual, and violence, with fine-grained subcategories such as hate/threatening, sexual/minors, and violence/graphic). Since the dataset is multi-label, we evaluate our model at the level of \emph{text–label pairs}: each annotation is treated as a separate prediction instance. Thus, while the corpus comprises 1,680 source texts, the effective number of evaluation instances is larger, reflecting the total set of text–label assignments. This framing allows a more precise assessment of category-specific moderation performance.  
(ii) a \textbf{Multimodal Moderation Dataset}~\cite{levi+al:25} of $1{,}500$ short-form videos spanning three categories: Death, Injury and Military Conflict; Drugs, Alcohol and Tobacco; and Kids content, across 12 languages and four modalities (text, thumbnail, transcript, video/frame). Together, these datasets enable evaluation of both text-only and multimodal moderation scenarios.  
\begin{table}[t]
\centering
\caption{\textbf{Summary of datasets.}}
\small
\setlength{\tabcolsep}{3.5pt}
\renewcommand{\arraystretch}{1.05}
\begin{tabularx}{\columnwidth}{@{}%
  >{\raggedright\arraybackslash}p{0.22\columnwidth}%
  >{\raggedright\arraybackslash}p{0.12\columnwidth}%
  >{\raggedright\arraybackslash}p{0.22\columnwidth}%
  >{\raggedright\arraybackslash}X%
@{}}
\toprule
\textbf{Dataset} & \textbf{Size} & \textbf{Languages / Modalities} & \textbf{Categories} \\
\midrule
OpenAI Moderation~\cite{markov2023holisticapproachundesiredcontent} &
1{,}680 texts &
English / Text &
Hate, Self-harm, Sexual, Violence (with subcategories) \\
\midrule
Multimodal Moderation~\cite{levi+al:25} &
1{,}500 videos &
12 languages / Text, Thumbnail, Transcript, Video/Frames &
DIMC (Death/Injury/Military Conflict), DAT (Drugs/Alcohol/Tobacco), Kids \\
\bottomrule
\end{tabularx}
\label{tab:datasets}
\end{table}
\subsection{Baselines}
We compare our LPP-based meta-model against widely used uncertainty heuristics:
\begin{itemize}
    \item \textbf{MSP}~\cite{hendrycks2017a,geifman2017selectiveclassificationdeepneural}.
    \item \textbf{Top-2 Margin}~\cite{gupta2024language}.
    \item \textbf{Entropy}~\cite{kendall+al:17,barshalom2025tokenprobabilitieslearnablefast}.
\end{itemize}
These represent the strongest post-hoc uncertainty signals. We also include a cost-insensitive \emph{always-trust} baseline.
\subsection{Evaluation Metrics}
To address \textbf{RQ1}, we evaluate error prediction using:
\begin{itemize}
    \item \textbf{AUC-ROC}: area under the receiver operating characteristic curve, measuring ranking quality of correct vs.\ incorrect predictions across thresholds.
    \item \textbf{F1-Score}: harmonic mean of precision and recall, balancing false positives and false negatives.
    \item \textbf{Macro-F1}: F1 macro-averaged over correctness labels (correct vs.\ incorrect), giving equal weight to error and non-error cases.
\end{itemize}
To address \textbf{RQ2}, we evaluate:
\begin{itemize}
    \item \textbf{Expected Cost $\mathcal{C}(\tau)$}, as defined in Sec.~\ref{subsec:policy}
    \item \textbf{Escalation Ratio}, defined as the fraction of items routed to human review:
    \[
    \text{Escalation Ratio} = \frac{TN + FN}{TP + FP + TN + FN}.
    \]
    Following best practice in selective prediction~\cite{geifman2017selectiveclassificationdeepneural,hendrycks2017a}, we report both the \emph{percentage} of items escalated and the \emph{absolute number of escalations}.
%    \item \textbf{Error Recall @ K}, simulating fixed human budgets ($K=10,20,45$)
\end{itemize}